%%%%%%%%%%%%%%%%%%%%%%%%%%%%%%%%%%%%%%%%%%%%%%%%%%%%%
%			HLAVIČKA								%
%%%%%%%%%%%%%%%%%%%%%%%%%%%%%%%%%%%%%%%%%%%%%%%%%%%%%
\documentclass[openany,12pt]{memoir}

\usepackage[utf8]{inputenc} 
\usepackage[czech]{babel}
\usepackage[T1]{fontenc}
\usepackage[top=1.5cm, bottom=2cm, left=2cm, right=2cm]{geometry}  % --> NASTAVENÍ OKRAJŮ
\usepackage{fancyhdr}
\usepackage{graphicx}
%\usepackage{lmodern}
\usepackage{tgschola}  % --- Silnější font
\usepackage{xwatermark}
\usepackage{xcolor}
\usepackage{changepage}
\usepackage{newunicodechar}



%%%%%% Package na zpěvník
\usepackage[full]{leadsheets}%http://mirrors.nic.cz/tex-archive/macros/latex/contrib/leadsheets/leadsheets_en.pdf   --> dokumentace	
\definesongtitletemplate{empty}{} 
\setchords{
format = \bfseries,   %tučné akordy
minor = {mi},% 
input-notation = {german},%
output-notation = {german}%
}
\definesongtitletemplate{empty}{} 

\newlength{\drop}
\newwatermark[allpages,color=red!50,angle=0,scale=2, xpos=0,ypos=0]{\includegraphics[width=5cm]{obr/pozadi2.jpg}} %--> dvojka na pozadí


%%%% Vlastní příkazy
\newcounter{Slokočet}   %Automatické číslování slok
\newcommand{\mezera}{\vspace*{0.5cm}}   %Horizontální odsazení slok
\newcommand{\stred}{5.2cm}   %%% Na zarovnání slok doprostřed, pozn. automatičtější zarovnávání na střed nejde
\newcommand{\refren}{\mezera \noindent \textbf{R:} } %refrén
\newcommand{\sloka}{\mezera \noindent \addtocounter{Slokočet}{1} \arabic{Slokočet}. } 	%sloka, která se automaticky čísluje
\newcommand{\ssloka}{\mezera \noindent}  % vlastní číslo sloky
\newcommand{\carka}{,\:}
\newcommand{\predehra}{\mezera \hspace*{-1.85cm} \textbf{Předehra:} }
\newcommand{\m}[1]{\color{white}{#1}}  %Pro akordy
\newcommand{\ap}{'} %Pro apostrof
\newcommand{\elipsa}{\kern\fontdimen3\font} %Příkaz pro lepší zacházení s výpustkami (=...); je to vpodstatě jen mezera mezi tečkama výpustky

%%%%%%%%%%%%%%%%%%%%%%%%%%%%%%%%%%%%%%%%%%%%%%%%%%%%%%
%			NÁVOD									 %
%%%%%%%%%%%%%%%%%%%%%%%%%%%%%%%%%%%%%%%%%%%%%%%%%%%%%%
%1. Věci v hlavičce IGNOROVAT
%2. Píseň psát do prostoru mezi \begin{song} a \end{song}
%3. další řádek se značí dvěma odsazeníma (= dvakrát stisknout enter)
%4. \refren vždy na začátku refrenu a \sloka na začátek sloky (automaticky se čísluje)
%  \ = alt gr + q ; [/] = alt gr f/g ; {/} = alt gr + b/n; ^ = alt gr + 3 + mezera
%Cokoliv napíšete do ^{  } se bude brát jako akord
%když se toto bude dotýkat nějakého slova (nebude mezi tím a slovem mezera)
%tak se akordy zjeví nad slovem, ale pište to před slovo
%Když se to nedotýká slova, tak akord lítá ve vzduchu a vytiskne se větší mezera
%První možnost je asi preferovanější
%5. Akordy stačí psát jen do první sloky, když se nezmění -- kytaristi to zvládnou



%\usepackage{subfiles}
%</++>

\begin{document}

\pagestyle{empty}
%\addtocontents{toc}{\protect\thispagestyle{empty}}
%\tableofcontents \thispagestyle{empty}\newpage
\newgeometry{top=1.5cm, bottom = 0cm, left = 2cm, right = 2cm}
	
\newpage

%Pro kompilaci je důležitý, 	aby ten input a \end{document} zůstali na posledních
%dvou řádcích!!!!


\begin{song}{title=\centering Bastila \\\normalsize Znouzecnost \vspace*{-0.3cm}}  %% sem se napíše jméno songu a autor
\moveright 4.5cm \vbox{      %Varianta č. 1  ---> Jeden sloupec zarovnaný na střed	

{\color{white}\_\_\_}


\sloka
^{Dmi{\color{white}\_}}Láska je nemoc ^{A#}a pravda přítěž

^{C{\color{white}\_\_\_}}Ideje maska ^{A}za kterou leccos zmizí

Slova jsou náboje pro ústa velké ráže

Peníze náboženství - jak cesta tak i cíl


\sloka
Srdce z olova ruce od špíny

Hovna ze zlata zaplatí všechny viny

Na cestě po žebříku vedoucím do nebe

Dobré je dolů kopat a smát se nad sebe

\refren
^{Dmi{\color{white}\_\_}}Bastila ^{A#{\color{white}\_\_}}padá ^{C}s ní hlavy ^{A{\color{white}\_\_\_}}králů

^{A#}z chaosu ^{C{\color{white}\_\_\_}}vzejde nový ^{A\,\,}řád

^{Dmi{\color{white}\_}}lůza se ^{A#{\color{white}\_}}baví ^{C}a ^{A}nevědomky

^{A#{\color{white}\_\_\_}}začíná ^{C{\color{white}\_\_\_\_}}stavět jiný ^{Dmi{\color{white}\_\_\_}}kriminál ^{A#\,\,C\,\,A}


\sloka
Bůh zemřel mlčky první den války

Zbyli je šamani a dobře ví jak na to

Nejvyšší umění - stvoření krávy z davů

A ta se krmí gesty a dojí moc a slávu


\sloka
"Všichni jsme bratři - ať žije svoboda"

- Jen jedna otázka - co komu kolik za to?

"Všichni jsme bratři - ať žije svoboda"

Občane okamžik - co ty mně - co já tobě?


\refren


\sloka
Zrozeni v kleci v kleci i umřem

A klec všech klecí nosíme ve své hlavě

Nejvyšší spravedlnost - smutný pán na orloji

Všichni jsme vyšli z prachu

Tam si budem zas rovni...


\refren


}
\setcounter{Slokočet}{0}
\end{song}

\end{document}
